% ============================================================================
% Fichier : annexes/Ann_C_templates.tex
% Titre   : Modèles de Documents Opérationnels
% Version : 1.0 -- novembre 2026
% ============================================================================
% Modèles annoncés dans le mode d'emploi (charte, RACI, SLA, grille de
% capacité), complétés de la demande de changement, de la fiche
% d'investissement, du constat en cinq éléments et de la chronologie.
% Contenus repris des chapitres 1, 3, 8, 11, 14, 15 et 16.
% ============================================================================

\chapter{Modèles de Documents Opérationnels}
\label{ann:templates}

Chaque modèle reprend les rubriques présentées dans un chapitre. Ils sont
à adapter à la taille de l'organisation~: un modèle complet mal tenu vaut
moins qu'un modèle simple tenu à jour.

\section*{1. Charte d'un comité stratégique des SI (Chapitre~\ref{chap:gouv-entreprise})}

\begin{tableaucours}{L{3.2cm}L{9.0cm}}{%
    \tableauHeader{Rubrique} &
    \tableauHeader{Contenu}
}
Mission & Assister le conseil dans l'évaluation de la stratégie numérique,
de la valeur des investissements et des risques liés au SI \\
Composition & Au moins deux administrateurs, dont le président~; DG et DSI
membres~; audit interne invité permanent \\
Pouvoirs & Accès direct à la DSI et à l'audit interne~; droit de demander
une expertise externe \\
Fréquence & Au moins trimestrielle~; sur convocation pour un projet majeur
en difficulté \\
Quorum et décisions & Règle de quorum~; décisions inscrites au procès-verbal
avec un responsable et une échéance \\
Reddition de comptes & Rapport écrit au conseil après chaque séance \\
\end{tableaucours}

\section*{2. Fiche d'investissement (Chapitre~\ref{chap:alignement})}

\begin{tableaucours}{L{3.2cm}L{9.0cm}}{%
    \tableauHeader{Rubrique} &
    \tableauHeader{Contenu}
}
Objectif & Objectif d'entreprise servi, et objectif d'alignement \\
Coût & Coût total de possession sur la durée de vie, par poste \\
Bénéfices & Bénéfices attendus, datés, avec leur indicateur et leur
propriétaire \\
Rentabilité & Valeur actuelle nette, taux d'actualisation retenu \\
Risques & Risques du projet et risques de ne pas le faire \\
Décision & Organe qui décide, date, conditions~; revue des bénéfices
prévue \\
\end{tableaucours}

\section*{3. Grille d'évaluation de la capacité d'un processus}

\begin{tableaucours}{L{2.8cm}cL{4.4cm}L{3.4cm}}{%
    \tableauHeader{Pratique} &
    \tableauHeader{Niv.} &
    \tableauHeader{Activité} &
    \tableauHeader{Preuve et cote}
}
BAI06.01 & 2 & \ldots & Pièce~; note~: complètement, largement,
partiellement, pas \\
BAI06.01 & 3 & \ldots & \ldots \\
\ldots & \ldots & \ldots & \ldots \\
\end{tableaucours}

\noindent Notes~: complètement (plus de 85~\%), largement (50 à 85~\%),
partiellement (15 à 50~\%), pas (moins de 15~\%). Un niveau est atteint
s'il est noté au moins « largement » et si tous les niveaux inférieurs sont
notés « complètement »~; la capacité retenue est celle que les preuves
établissent, non celle que l'on déclare (Chapitre~\ref{chap:cobit-capacite}).

\section*{4. Demande de changement (Chapitre~\ref{chap:itil-changements})}

\begin{tableaucours}{L{3.2cm}L{9.0cm}}{%
    \tableauHeader{Rubrique} &
    \tableauHeader{Contenu}
}
Description & Ce qui change, pourquoi, éléments de configuration concernés \\
Type & Standard, normal ou urgent \\
Risque & Impact et probabilité~; services et clients exposés (analyse
d'impact sur la CMDB) \\
Tests & Recette réalisée, taux, anomalies bloquantes closes \\
Mise en production & Date dans le calendrier, période de gel, équipes
mobilisées \\
Retour arrière & Plan testé, critères de déclenchement, heure limite \\
Autorisation & Autorité de changement compétente, décision motivée \\
Revue & Résultat constaté après la mise en production \\
\end{tableaucours}

\section*{5. Accord de niveau de service (Chapitre~\ref{chap:itil-niveaux})}

\begin{tableaucours}{L{3.2cm}L{9.0cm}}{%
    \tableauHeader{Rubrique} &
    \tableauHeader{Contenu}
}
Description du service & Ce que le service permet au client d'obtenir \\
Horaires & Périodes où les engagements s'appliquent \\
Disponibilité & Part du temps de service convenu où le service est
utilisable \\
Délais de rétablissement & Par priorité \\
Indicateur de résultat & Ce que le client vit, défini avec lui \\
Mesure et compte rendu & Source des données, fréquence \\
Pénalités & Conséquences du non-respect, et leur application \\
Revue & Fréquence de révision, comité de service \\
\end{tableaucours}

\section*{6. Matrice RACI}

\begin{tableaucours}{L{4.0cm}ccccc}{%
    \tableauHeader{Activité} &
    \tableauHeader{Rôle 1} & \tableauHeader{Rôle 2} &
    \tableauHeader{Rôle 3} & \tableauHeader{Rôle 4} &
    \tableauHeader{Rôle 5}
}
Activité 1 & A & R & C & I & \\
Activité 2 & & A & R & C & I \\
\ldots & & & & & \\
\end{tableaucours}

\noindent Règles (Chapitre~\ref{chap:processus})~: un seul A par activité~;
au moins un R~; le A au bon niveau~; peu de C~; des fonctions et non des
personnes.

\section*{7. Constat d'audit en cinq éléments (Chapitre~\ref{chap:audit})}

\begin{tableaucours}{L{3.2cm}L{9.0cm}}{%
    \tableauHeader{Élément} &
    \tableauHeader{Question}
}
Constatation & Que se passe-t-il réellement~? \\
Critère & Que devrait-il se passer~? Avec sa source \\
Cause & Pourquoi l'écart existe-t-il~? \\
Conséquence & Quel risque, ou quel effet~? \\
Recommandation & Que faut-il faire, par qui, pour quand~? \\
\end{tableaucours}

\section*{8. Chronologie de défaillance}

\begin{tableaucours}{L{2.6cm}L{6.4cm}L{3.2cm}}{%
    \tableauHeader{Date} &
    \tableauHeader{Fait établi} &
    \tableauHeader{Pièce}
}
\ldots & Fait, décision, ou absence de décision & Procès-verbal, courriel,
journal \\
\end{tableaucours}

\noindent La chronologie ne contient que des constatations~; l'analyse vient
ensuite, avec son niveau de confiance (Chapitre~\ref{chap:prologue}).
